\documentclass{article}

\usepackage{comment, multicol}
\usepackage{hyperref}

\usepackage{calc,pict2e,picture}
\usepackage{textgreek, textcomp, gensymb, stix}

\setcounter{secnumdepth}{2}

\title{\LaTeX.js Showcase}
\author{made with $\varheartsuit$ by Michael Brade}
\date{2017--2021}

\begin{document}
\maketitle

\begin{abstract}
Hello This document will show most of the features of \LaTeX.js while at the same time being a
gentle introduction to \LaTeX. In the appendix, the API as well as the format of custom macro definitions
in JavaScript will be explained.
\end{abstract}

\section{Characters}

It is possible to input any UTF-8 character either directly or by character code using one
of the following:

\begin{itemize}
  \item \texttt{\textbackslash symbol\{"00A9\}}: \symbol{"00A9}
  \item \verb|\char"A9|: \char"A9
  \item \verb|^^A9 or ^^^^00A9|: ^^A9 or ^^^^00A9
\end{itemize}

\bigskip
\noindent
Special characters, like those:

\begin{center}
\$ \& \% \# \_ \{ \} \~{} \^{} \textbackslash
\end{center}

have to be escaped.

More than 200 symbols are accessible through macros. For instance: 30\,°C is 86\,°F.
See also Section~\ref{sec:symbols}.

\section{Spaces and Comments}

Spaces and comments, of course, work just as they do in \LaTeX. This is an example: Su-
percalifragilisticexpialidocious

It does not matter whether you enter one or several spaces after a word, it will always
be typeset as one space—unless you force several spaces, like\quad now.

New \TeX users may miss whitespaces after a command. Experienced \TeX users are
\TeX{}perts, and know how to use whitespaces.

Longer comments can be embedded in the \texttt{comment} environment: This is another ex-
ample for embedding comments in your document.

\section{Dashes and Hyphens}

\LaTeX knows four kinds of dashes. Access three of them with different numbers of consecu-
tive dashes. The fourth sign is actually not a dash at all—it is the mathematical minus
sign:

\begin{center}
daughter-in-law, X-rated\\
pages 13--67\\
yes---or no? \\
$0$, $1$ and $-1$
\end{center}

The names for these dashes are: ‘-’ hyphen, ‘--’ en-dash, ‘---’ em-dash,
and ‘$-$’ minus sign.

\end{document}
